% 報告会テンプレート（LuaLaTeX対応版）
\documentclass[10pt,a4j,twocolumn]{ltjsarticle}


\usepackage{graphicx}   % \includegraphics を使うためのパッケージ読み込み
\usepackage{imcreport}  % IMC報告会用テンプレートパッケージの読み込み


% ハイフネーションを抑制する．数値が大きいほど，無理にでも抑制しようとする．
\hyphenpenalty=10000\relax
\exhyphenpenalty=10000\relax
\sloppy


\title{中間報告用Lua\LaTeX テンプレートの使い方} % タイトル
\author{発表者氏名}                              % 著者
\studentnumber{学籍番号}                         % 学籍番号
\date{YYYY--MM--DD}                              % 日付
\header{中間報告}                                % 左上のヘッダ内容


\begin{document}

\maketitle % タイトル，著者情報などを挿入

\section{はじめに}

報告会で使われるフォーマットに従ったLua\LaTeX のスタイル，\verb|imcreport.sty|を作りました．この資料は作成したスタイルを使ったテンプレートであり，ここにはテンプレートの使い方が書かれています．

作成したスタイルはLua\LaTeX から利用され，クラスは\verb|ltjsarticle|が指定されてていることを前提にしています．
つまりこのスタイルは\pLaTeX, u\pLaTeX であったり，クラスに\verb|ltjsarticle|以外が指定されてる場合の動作を保証しません．

以下ではテンプレートのソースコード，つまり\verb|main.tex|から抜粋してテンプレートの使い方を説明することがあります．
したがってこのテンプレートのPDFとソースコードの両方を見比べながら読んでください．

\section{図の挿入}

Lua\LaTeX では図~\ref{バッテンのラベル}のように，図を挿入することができます．
\begin{figure}[bp]
  \centering                                                  % 画像とキャプションを中央揃えにする
  \includegraphics[width=80truemm]{./sample_figure.pdf}       % 指定幅（width）で画像（pdf, png, jpg, eps）を挿入．
                                                              % 理由は省略するが，寸法にmmではなくtruemmを使うこと．
  \caption{バッテン}                                          % キャプション
  \label{バッテンのラベル}                                    % 本文から参照するためのラベル
\end{figure}
\verb|\begin{figure}|の後ろに位置を指示する文字を四角カッコ\verb|[]|を使って与えることで画像の挿入位置を指定することができます．
\verb|t|, \verb|b|, \verb|p|, \verb|h|の4種類を指定することができ，それぞれ次のような意味を持ちます．
\begin{description}
  \item[t] ページ上部に挿入
  \item[b] ページ下部に挿入
  \item[p] 新規ページとして挿入
  \item[h] 現在位置に挿入
\end{description}
ただし必ずしも指示どおりにはならず，Lua\LaTeX エンジンが見栄えのいいところに自動で移動させる場合があります．
勝手に画像を動かしてしまうと考えるか，適当でもそれなりの位置に挿入してくれると考えるかは難しいところです……．
位置を指定しない場合は\verb|[tbp]|として処理されます．

\section{表の挿入}

Lua\LaTeX で作った表を表~\ref{表のラベル}に示します．
こちらも図と同様の方法で挿入位置を指定することができます．
勝手に動かされるのも同じです．
\begin{table}
  \centering
  \caption{表のタイトル}
  \label{表のラベル}
  \begin{tabular}{c|cccc}
    \hline
          & 項目A & 項目B & 項目C & 項目D \\
    \hline \hline
    項目a & 1     & 2     & 3     & 4     \\
    項目b & 5     & 6     & 7     & 8     \\
    \hline
  \end{tabular}
\end{table}

\verb|\begin{tabular}{XXX}|の\verb|XXX|部分で列の整列方向と垂直線の位置を指定します．
表~\ref{表のラベル}の場合は\verb|XXX|部分は\verb=c|cccc=になっています．
これは表が``中央揃え''，``垂直線''，``中央揃え''，``中央揃え''，``中央揃え''，``中央揃え''の6列と1つの垂直線で構成されることを意味します．
\verb|c|が中央揃えの列，\verb|l|（エル）が左揃えの列，\verb|r|が右揃えの列であることを意味します．
列と違って行はあらかじめ行や垂直線の配置を指定する必要はなく，書いた分だけ勝手に縦に伸びます．
行の区切りは\verb|\\|で，列の区切りは\verb|&|です．
水平線は\verb|\hline|で引けます．
より具体的な書き方はソースコードを参照してください．

\section{セクションとサブセクション}

報告書では話題を分けるために見出しを使います．
主にセクション（章または節）とサブセクション（節または項）の2種類を使うことになるでしょう．

\subsection{セクション（section）}

\verb|\section{セクション名}|で作れます．
このテンプレートでは``はじめに''，``図の挿入''，``表の挿入''，``セクションとサブセクション''がセクションです．

\subsection{サブセクション（subsection）}

\verb|\subsection{サブセクション名}|で作れます．
このテンプレートでは``セクション（section）''，``サブセクション（subsection）''がサブセクションです．

\end{document}

